\documentclass[a4paper,11pt]{ltjsarticle}
\usepackage{geometry}
\geometry{top=25mm,bottom=25mm,left=25mm,right=25mm}
\usepackage{listings}
\usepackage{booktabs}
\usepackage{fancyhdr}
\usepackage{hyperref}
\usepackage{xcolor}
\usepackage{amsmath}
\usepackage{amssymb}
\usepackage{longtable}
\usepackage{array}

\lstset{
  basicstyle=\ttfamily\small,
  keywordstyle=\color{blue}\bfseries,
  commentstyle=\color{green!50!black},
  stringstyle=\color{red},
  numbers=left, numberstyle=\tiny\color{gray}, numbersep=5pt,
  frame=single, breaklines=true, tabsize=4, showstringspaces=false, captionpos=b
}

\pagestyle{fancy}
\fancyhf{}
\fancyhead[L]{計算機科学実験及演習3 ハードウェア — 実施報告書}
\fancyfoot[C]{\thepage}
\renewcommand{\headrulewidth}{0.4pt}

\begin{document}

\begin{titlepage}
\centering
\vspace*{3cm}
{\LARGE 計算機科学実験及演習3 ハードウェア}\\[1cm]
{\Huge\bfseries 実施報告書}\\[0.5cm]
{\LARGE SIMPLE/B 拡張プロセッサの設計・実装}\\[3cm]
\begin{tabular}{rl}
学籍番号 & （学籍番号を記入） \\[0.5em]
氏名     & （氏名を記入） \\[0.5em]
提出日   & 2026年6月5日 \\
\end{tabular}
\vfill
京都大学 工学部情報学科 計算機科学コース
\end{titlepage}

\tableofcontents
\clearpage


% =====================================================================
\section{プロジェクト概要}

本プロジェクトでは、SIMPLE アーキテクチャ（ver 4.0）に基づく16ビットプロセッサを
Verilog HDL で設計・実装し、PowerMedusa MU500-RX/RK (Cyclone IV FPGA) 上で動作させた。
基本のSIMPLE/Bプロセッサに対し、マイクロアーキテクチャ改良、ISA拡張、
周辺回路追加、応用プログラム開発を行った。


% =====================================================================
\section{役割分担}

2名のグループで以下の通り作業を分担した。

\begin{table}[h]
\centering
\caption{役割分担}
\begin{tabular}{clp{8cm}}
\toprule
担当 & 主要範囲 & 具体的な作業 \\
\midrule
A & CPUコア設計 &
  4サイクル実行方式の設計・実装、
  ISA拡張（BAL/BR/ADDI）、
  同期RAM（Block RAM）対応、
  テストプログラム（test\_all.hex）の作成 \\
\midrule
B & 周辺回路・ソフトウェア &
  ハードウェアタイマ（timer.v）、
  I/Oアドレス空間拡張、
  8桁7セグ表示+モード切替、
  バブルソートプログラム（sort.asm）、
  Pythonアセンブラ（simple\_asm.py）、
  テストベンチ拡張 \\
\bottomrule
\end{tabular}
\end{table}

\subsection{共通作業}

\begin{itemize}
\item 基本設計仕様の策定（plan.md）
\item FPGA実装とデモ準備
\item 最終レポート群の作成
\end{itemize}


% =====================================================================
\section{実施スケジュール}

\begin{table}[h]
\centering
\caption{実施スケジュール}
\begin{tabular}{lp{5.5cm}p{5.5cm}}
\toprule
期間 & 担当A & 担当B \\
\midrule
4月第3週 &
  SIMPLE/B基本設計の理解、5フェーズCPU実装 &
  導入課題（7セグ、カウンタ）完了 \\
\midrule
4月第4週 &
  4サイクル化の設計、パイプラインレジスタ導入 &
  SIMPLE/B基本設計の理解、plan.md策定 \\
\midrule
5月第1週 &
  ISA拡張（BAL/BR/ADDI）実装、テストプログラム作成 &
  タイマモジュール設計・実装 \\
\midrule
5月第2週 &
  同期RAM対応、機能設計仕様書（担当A）作成 &
  I/Oアドレス拡張、7セグ8桁化 \\
\midrule
5月第3--4週 &
  統合テスト支援、中間レポート &
  バブルソートプログラム、アセンブラ、中間レポート \\
\midrule
5月第5週--6月 &
  性能評価、FPGA最終調整 &
  テストベンチ拡張、最終レポート \\
\bottomrule
\end{tabular}
\end{table}


% =====================================================================
\section{実装結果}

\subsection{実装状況一覧}

\begin{table}[h]
\centering
\caption{実装項目の状況}
\begin{tabular}{clcc}
\toprule
\# & 項目 & 計画 & 達成 \\
\midrule
1 & 4サイクル実行方式         & $\checkmark$ & $\checkmark$ \\
2 & ISA拡張 (BAL)             & $\checkmark$ & $\checkmark$ \\
3 & ISA拡張 (BR)              & $\checkmark$ & $\checkmark$ \\
4 & ISA拡張 (ADDI)            & $\checkmark$ & $\checkmark$ \\
5 & 同期RAM (Block RAM) 対応  & $\checkmark$ & $\checkmark$ \\
6 & ハードウェアタイマ         & $\checkmark$ & $\checkmark$ \\
7 & I/Oアドレス空間 (6ポート) & $\checkmark$ & $\checkmark$ \\
8 & 8桁7セグ表示              & $\checkmark$ & $\checkmark$ \\
9 & 表示モード切替 (4モード)  & $\checkmark$ & $\checkmark$ \\
10& LED出力                    & $\checkmark$ & $\checkmark$ \\
11& バブルソートプログラム     & $\checkmark$ & $\checkmark$ \\
12& Pythonアセンブラ           & $\checkmark$ & $\checkmark$ \\
\bottomrule
\end{tabular}
\end{table}

計画した全12項目を完了した。

\subsection{ファイル構成}

最終的なファイル構成を表\ref{tab:files}に示す。

\begin{table}[h]
\centering
\caption{最終ファイル構成}
\label{tab:files}
\small
\begin{tabular}{lll}
\toprule
ファイル & 種別 & 概要 \\
\midrule
\texttt{simple\_top.v}   & Verilog & トップモジュール（拡張版） \\
\texttt{simple\_cpu.v}   & Verilog & CPUコア（4サイクル+ISA拡張） \\
\texttt{alu.v}           & Verilog & ALU \\
\texttt{shifter.v}       & Verilog & バレルシフタ \\
\texttt{regfile.v}       & Verilog & レジスタファイル \\
\texttt{ram.v}           & Verilog & Block RAM (Quartus IP) \\
\texttt{ram\_sim.v}       & Verilog & シミュレーション用RAM \\
\texttt{seg7dec.v}       & Verilog & 7セグデコーダ \\
\texttt{timer.v}         & Verilog & ハードウェアタイマ \\
\texttt{simple\_tb.v}     & Verilog & テストベンチ \\
\midrule
\texttt{simple\_asm.py}   & Python  & アセンブラ \\
\texttt{sort.asm}        & Assembly & バブルソートプログラム \\
\texttt{sort.hex}        & データ  & バブルソート機械語 \\
\texttt{test\_all.hex}    & データ  & ISA全命令テスト \\
\texttt{test\_basic.hex}  & データ  & 基本テスト \\
\texttt{test\_branch.hex} & データ  & 分岐テスト \\
\bottomrule
\end{tabular}
\end{table}


% =====================================================================
\section{検証結果}

\subsection{シミュレーション検証}

Icarus Verilog 11.0 による自動検証を行い、全テストがパスした。

\subsubsection{ISA全命令テスト}

テストプログラム \texttt{test\_all.hex} により、基本命令および拡張命令の
正常動作を5項目で検証した。

\begin{table}[h]
\centering
\caption{ISAテスト結果}
\begin{tabular}{clcc}
\toprule
\# & テスト内容 & 期待出力 & 結果 \\
\midrule
0 & 即値ロード + 加算 (LI, ADD)          & 0x0008 & OK \\
1 & ループ (ADD, CMP, BLT)               & 0x0003 & OK \\
2 & 即値加算 (LI, ADDI)                  & 0x000A & OK \\
3 & サブルーチン呼出 (BAL, BR)            & 0x0014 & OK \\
4 & サブルーチン内レジスタ変更 (LI)       & 0xFFFF & OK \\
\bottomrule
\end{tabular}
\end{table}

\subsubsection{バブルソートテスト}

テストプログラム \texttt{sort.hex}（58語）により、
16要素の配列ソートと結果出力を検証した。

\begin{lstlisting}[language={},frame=single,caption={バブルソートテスト結果}]
=== Bubble Sort Test ===
OUT[0]  = 0000  OK    OUT[8]  = 0008  OK
OUT[1]  = 0001  OK    OUT[9]  = 0009  OK
OUT[2]  = 0002  OK    OUT[10] = 000a  OK
OUT[3]  = 0003  OK    OUT[11] = 000b  OK
OUT[4]  = 0004  OK    OUT[12] = 000c  OK
OUT[5]  = 0005  OK    OUT[13] = 000d  OK
OUT[6]  = 0006  OK    OUT[14] = 000e  OK
OUT[7]  = 0007  OK    OUT[15] = 000f  OK
=== CPU Halted (PC=003a, 16 outputs) ===
=== SORT TEST PASSED ===
\end{lstlisting}

初期データ \{15, 3, 8, 1, 12, 5, 10, 2, 14, 7, 11, 4, 13, 6, 9, 0\} が
昇順 \{0, 1, 2, ..., 15\} に正しくソートされた。

\subsection{FPGA実機検証}

（実機デモの結果を記載すること）


% =====================================================================
\section{開発中に遭遇した課題と解決}

\subsection{同期RAMのタイミング問題}

\textbf{課題}: Block RAM は非同期リードではなく、クロックエッジでアドレスをラッチする。
基本設計のメモリアクセスタイミングでは、フェッチしたアドレスの命令が同じサイクルで得られず、
正しく動作しなかった。

\textbf{解決}: データが必要な1フェーズ前にアドレスを提示する方式を採用。
PH3でLD/ST用アドレス、PH4で次命令アドレスを提示することで、
4サイクル実行を維持したまま同期RAMと協調動作させた。

\subsection{4サイクル化でのレジスタ書き戻し}

\textbf{課題}: PH1でフェッチと書き戻しを同時に行うため、
前命令の演算結果（DR/MDR）の選択情報をPH4$\to$PH1間で保持する必要があった。

\textbf{解決}: パイプラインレジスタ（\texttt{wb\_rf\_we}, \texttt{wb\_rf\_addr},
\texttt{wb\_rf\_data}, \texttt{wb\_use\_mdr}）を導入。
PH4で書き戻し情報を保存し、PH1で使用する。
レジスタファイルの非同期リードにより、PH2で最新値が読み出されるため
データハザードは発生しない。

\subsection{バブルソートの外側ループ終了条件}

\textbf{課題}: 初期実装では外側ループの終了判定に CMP + BNE を使用していたが、
i=0を超えてマイナスになっても BNE が成立し続け、無限ループとなった。

\textbf{解決}: ADDI命令がフラグを更新する性質を利用し、
\texttt{ADDI r3, -1} のZフラグを直接 \texttt{BNE} で参照する方式に変更。
i=1$\to$0 のとき Z=1 となり BNE 不成立でループ脱出。
CMP命令が不要になり、1命令削減にもなった。

\subsection{I/Oポートの後方互換性}

\textbf{課題}: I/Oポート番号を導入すると、ポート番号を指定しない
既存プログラムが動作しなくなる可能性があった。

\textbf{解決}: IN/OUT命令の\texttt{d[3:0]}フィールドが0の場合にポート0を選択する設計とした。
既存プログラムではこのフィールドは0であるため、そのまま動作する。


% =====================================================================
\section{まとめと今後の課題}

\subsection{達成事項}

\begin{itemize}
\item 計画した全12項目を完了し、機能検証をパスした
\item 4サイクル化により25\%の性能向上を達成
\item ISA拡張によりサブルーチン呼出と即値演算を実現
\item タイマとI/Oポートにより周辺デバイス制御が可能になった
\item 58語のバブルソートプログラムが正常に動作
\item アセンブラによりプログラム開発の効率が大幅に向上
\end{itemize}

\subsection{今後の課題}

\begin{itemize}
\item \textbf{パイプライン化}: 4サイクルから更に進んで命令パイプラインを導入すれば、
  理想的にはCPI=1（4倍の性能向上）が期待できるが、
  データハザード・コントロールハザードの対策が必要
\item \textbf{割込み}: タイマのoverflowを割込みとして処理できれば、
  ポーリング不要で効率的なタイマ利用が可能になる
\item \textbf{乗算器}: 16ビット乗算命令の追加により、
  より複雑な数値計算プログラムの実現が容易になる
\item \textbf{キャッシュ}: メモリ階層を導入し、
  Block RAMをキャッシュとして外部メモリを接続する構成
\end{itemize}


\end{document}
